\pdfobjcompresslevel=0
\documentclass[11pt,a4paper,oneside]{book}
\usepackage[T1]{fontenc}
\usepackage[utf8]{inputenc}
\usepackage{lmodern}
\input{glyphtounicode}
\pdfgentounicode=1
% Large integral glyphs otherwise extract as Latin Z in TeX Live 2022.
\pdfglyphtounicode{integraldisplay}{222B}
\pdfglyphtounicode{integraltext}{222B}
\usepackage[ngerman]{babel}
\usepackage{amsmath,amssymb,graphicx,booktabs,longtable,listings,xcolor,makeidx}
\usepackage[margin=28mm]{geometry}
\definecolor{navy}{HTML}{0B1F2A}
\definecolor{accent}{HTML}{168557}
\usepackage[unicode,colorlinks=true,linkcolor=navy,urlcolor=accent,citecolor=navy]{hyperref}
% Permit line breaks in long repository paths without another TeX dependency.
\expandafter\def\expandafter\UrlBreaks\expandafter{\UrlBreaks
\do\a\do\b\do\c\do\d\do\e\do\f\do\g\do\h\do\i\do\j\do\k\do\l\do\m
\do\n\do\o\do\p\do\q\do\r\do\s\do\t\do\u\do\v\do\w\do\x\do\y\do\z
\do\A\do\B\do\C\do\D\do\E\do\F\do\G\do\H\do\I\do\J\do\K\do\L\do\M
\do\N\do\O\do\P\do\Q\do\R\do\S\do\T\do\U\do\V\do\W\do\X\do\Y\do\Z}
\input{metadata}
\hypersetup{pdftitle={FEM Toolbox -- Handbuch},pdfauthor={FEM Toolbox},pdfsubject={Version \ToolboxVersion; Commit \SourceCommit; KANDIDAT}}
\lstset{basicstyle=\small\ttfamily,breaklines=true,frame=single,columns=fullflexible}
\lstset{literate={²}{{$^2$}}1 {³}{{$^3$}}1 {µ}{{$\mu$}}1 {Δ}{{$\Delta$}}1 {α}{{$\alpha$}}1 {ε}{{$\varepsilon$}}1 {κ}{{$\kappa$}}1 {ᵀ}{{$^T$}}1}
\makeatletter
\renewcommand*\l@subsection{\@dottedtocline{2}{3.8em}{5.2em}}
\makeatother
\makeindex
\newif\ifBookFigures
\newif\ifBookTables
\AtBeginDocument{\renewcommand{\indexname}{Sachregister}}
\newcommand{\BookTitle}{Gesamthandbuch}
\newcommand{\FrontMatter}{%
\frontmatter
\begin{titlepage}
{\Huge\sffamily\bfseries FEM Toolbox\par}\vspace{12mm}
{\LARGE\BookTitle\par}\vspace{8mm}
{\Large KANDIDAT -- Gesamtabnahme\par}\vfill
Version \ToolboxVersion\\[4mm]
Quellcommit: {\small\ttfamily\SourceCommit}\\[4mm]
Quellzustand: \SourceState\\[4mm]
M9 / Issues 206--210. Vier Fachbände; redaktionell geprüfter Kandidat.\\
Keine experimentelle Validierungsfreigabe.\\[8mm]
Deutschsprachige Referenz mit englischen API-Bezeichnern.
\end{titlepage}
\tableofcontents
\ifBookFigures\listoffigures\fi
\ifBookTables\listoftables\fi
\chapter{Lesehinweise und Status}
Die vier Fachbände behandeln Theorie, Anwendung, API/Entwicklung und
Verifikation bis zum generischen M9-Unsicherheitsworkflow. Die Einzelabnahmen
206--210 sind integriert. Die redaktionelle Gesamtabnahme 211 wird im
Quellenarchiv protokolliert; Veröffentlichung und finale Versionsmetadaten
folgen mit Release-Issue 66. Der Kandidat bezeichnet den angegebenen Commit,
nicht eine neue Veröffentlichung der auf der Titelseite genannten Paketversion.

Experimentelle Validierung bleibt außerhalb M9 offen. Spätere Adaptivität,
Optimierung, RANS/FSI, SPICE und die Femantis-Namensmigration gehören nicht
zu diesem Funktionsstand. Historische Releaseberichte und Entwicklungspläne
werden als Herkunft behandelt. Es liegt kein unabhängiges Fachreview vor;
CI und eigene Quellen-/Layoutprüfung sind im Abnahmebericht getrennt benannt.
\section{Symbole und Abkürzungen}
\begin{tabular}{ll}
$\Omega$ & Gebiet in drei Dimensionen\\
$\mathbf n$ & äußere Einheitsnormale\\
$T$ & absolute Temperatur, K\\
$\mathbf q$ & Wärmestromdichte, W/m$^2$\\
FEM & Finite-Elemente-Methode\\
SI & Internationales Einheitensystem\\
V\&V & Verifikation und Validierung
\end{tabular}
\mainmatter}
\newcommand{\BackMatter}{%
\backmatter
\bibliographystyle{plain}\bibliography{references}
\printindex}
